\documentclass[12pt]{article}
\usepackage[utf8]{inputenc}
\usepackage[T1]{fontenc}
\usepackage{hyperref}
\title{Wiyot Language Resources for K-12 Curriculum Development}
\author{Compiled by Claude Sonnet 5.5 (claude-sonnet-5-5) for Andr\'as Kornai}
\date{}
\begin{document}
\maketitle

\section{Introduction}
Wiyot (ISO 639-3 wiy) is an Algic language of the Humboldt Bay area of northwestern California. It has no fluent native speakers. Della Prince, the last known native speaker, died in 1962. The federally recognized Wiyot Tribe runs a revitalization program, Soulatluk (the Wiyot ancestral language). Overview: \url{https://en.wikipedia.org/wiki/Wiyot_language}.

\section{Primary Educational Resources}
According to the search results, the Tribe offers language courses and publishes Wiyot texts, including an annual calendar. I did not see links to specific curricula or school programmes.

\section{Linguistic and Reference Materials}
A UC Berkeley linguistics project, started in 2014 with the Tribe's Language Program, is building a searchable dictionary and text corpus: \url{https://linguistics.berkeley.edu/~wiyot/about.html}. Its notes describe the dictionary as pre-beta. The site directs grammar and usage questions to the Tribe's Language Director. An NSF-funded documentation grant record is at \url{https://grantome.com/grant/NSF/BCS-0854176}. Language overview with sources: \url{https://languagegeek.com/california/wiyot.html}.

\section{Religious and Community Resources}
The ``Speaking Soulatluk'' project placed signs on the Eureka waterfront with QR codes linking to Wiyot language resources and audio: \url{https://www.inkpeople.org/speakingsoulatluk}. A news report on a federal grant for Native language revitalization that covers the Wiyot is at \url{https://krcb.org/2022061779834/news-feed/fading-native-languages-thrown-federal-lifeline}.

\section{Digital and Audio Resources}
The QR-linked audio recordings from the Eureka signs are the main digital resource I found. I did not verify what the Tribe's own website offers, and it is worth checking directly with the Tribe's Language Program.

\section{Curriculum Development Context}
Wiyot is a revitalization case with no first-language speakers, but it has an active tribal language program, a developing dictionary and corpus, and archival recordings. A K-12 heritage curriculum should be built in partnership with the Wiyot Tribe's Language Program.
\end{document}